% !TeX root = ../抽象代数笔记.tex
% 由主文档通过 \input 载入；请编译主文档。
\setcounter{chapter}{14} % 保留原书章号。
\chapter{尺规作图与二次扩张}\label{ch:constructions}

\begin{description}
  \item[学习目标] 将几何作图步骤转化为域运算，理解代数条件如何给出不可作图证明。
  \item[先修知识] 第\ref{ch:extensions}章，平面解析几何。本章为应用专题，可独立选学。
\end{description}

\section{可作点与可作数}
在欧氏平面中固定原点与单位长度，明确无刻度直尺和圆规允许的操作。定义可作点、可作实数（constructible real number），说明坐标系与初始数据的作用。先建立加、减、乘、除及对正数开平方的几何实现。

\section{交点方程与二次扩张塔}
分析直线与直线、直线与圆、圆与圆的交点方程，证明一次作图至多添加二次方程的根。建立实数 $\alpha$ 从有理初始数据可作，当且仅当它属于一条有限实域塔的顶层，其中起点为 $\mathbb Q$、每步扩张次数为 $1$ 或 $2$。分别证明必要性与充分性。

\section{次数障碍与经典问题}
由塔式公式推出可作代数数 $\alpha$ 的次数 $[\mathbb Q(\alpha):\mathbb Q]$ 必为 $2$ 的幂，并明确此条件本身不充分。用 $\sqrt[3]{2}$ 证明倍立方不可作，用三倍角关系与不可约三次多项式证明一般三等分角不可作；某些特定角仍可三等分。

\section{正多边形与论证边界}
将正多边形作图与单位根联系，正五边形作为核心例题。一般正多边形的可作判别在圆分扩张之后列为选学回访。化圆为方的讨论明确依赖 $\pi$ 的超越性；该超越性作为外部定理说明，不声称可由本章初等域论证明。

\section{例题、练习与自检}
构造若干含平方根的长度，判断给定代数数是否能由次数障碍排除。练习要求区分“不能由当前必要条件判定”与“已经证明可作”，并解释为什么证明某个角不可三等分足以否定普遍作图程序，却不能否定所有特殊角。

